\documentclass{article}
\usepackage[utf8]{inputenc}
\usepackage{amsmath, amssymb, amsthm}
\usepackage{tikz-cd}
\usepackage{hyperref}
\usepackage{listings}
\usepackage{xcolor}

\DeclareUnicodeCharacter{03C6}{\ensuremath{\varphi}}
\DeclareUnicodeCharacter{03A6}{\ensuremath{\Phi}}
\DeclareUnicodeCharacter{03B3}{\ensuremath{\gamma}}
\DeclareUnicodeCharacter{03BC}{\ensuremath{\mu}}
\DeclareUnicodeCharacter{03B7}{\ensuremath{\eta}}
\DeclareUnicodeCharacter{03C4}{\ensuremath{\tau}}
\DeclareUnicodeCharacter{03BA}{\ensuremath{\kappa}}
\DeclareUnicodeCharacter{2243}{\ensuremath{\simeq}}

\newtheorem{theorem}{Theorem}[section]
\newtheorem{lemma}[theorem]{Lemma}
\newtheorem{definition}[theorem]{Definition}
\newtheorem{remark}[theorem]{Remark}

\lstdefinelanguage{lean}{
    basicstyle=\ttfamily\small,
    keywordstyle=\color{blue}\bfseries,
    commentstyle=\color{gray},
    morekeywords={import, namespace, open, universe, abbrev, class, extends, where, instance, section, variable, theorem, def, by, intros, rfl, ext, dsimp, erw, rw, change, congr, exact, have, fun, let},
    morecomment=[l]{--},
    morecomment=[s]{/-}{-/},
    literate=
        {⥤}{{$\rightsquigarrow$}}1
        {⋙}{{$\ggg$}}1
        {⟶}{{$\longrightarrow$}}1
        {𝟭}{{$\mathbf{1}$}}1
        {𝟙}{{$\text{id}$}}1
        {γμ}{{$\gamma\mu$}}2
        {γη}{{$\gamma\eta$}}2
        {γ}{{$\gamma$}}1
        {μ}{{$\mu$}}1
        {η}{{$\eta$}}1
        {τ}{{$\tau$}}1
        {φ}{{$\varphi$}}1
        {Φ}{{$\Phi$}}1
        {κ}{{$\kappa$}}1
        {↾}{{$\upharpoonright$}}1
        {∀}{{$\forall$}}1
        {→}{{$\to$}}1
        {=>}{{$\Rightarrow$}}1
        {≫}{{$\gg$}}1
        {≃}{{$\simeq$}}1
}
\lstset{
    language=lean,
    breaklines=true,
    frame=single,
    showstringspaces=false
}

\title{Pointfree CoEndoYoneda Lemma \\ 
in \\
Functional Categories \\ 
Theory and Formalization in Lean 4}
\author{Luc Duponcheel \and Antigravity AI}
\date{\today}

\begin{document}

\maketitle

\begin{abstract}
We present a variation of the CoYoneda Lemma for functional categories equipped with a monadic
structure of its global elements, which we formally identify as \emph{Functional Categories}. The
central theoretical contribution, the Pointfree CoEndoYoneda Equivalence, establishes an isomorphism
between, on the one hand, natural transformations mapping out of co-Yoneda endofunctors to a
designated kind of global endofunctors and, on the other hand, some corresponding kind of global
elements. 

The original motivation was to formulate an equivalence that could play a role in the broader
programming  world, encompassing both pure programs, more precisely functions, and impure programs,
more precisely computation with side-effects valued functions (modeled via Kleisli categories).
While the requirements proved too restrictive to admit arbitrary monadic effects, this limitation
reveals a profound  mathematical result: the \emph{Functional Category} constraints categorically
classify the essence of referential transparency.

In addition to outlining the mathematical foundations, we detail the complete formalization of this
result in the Lean 4 theorem prover. We discuss the technical challenges of working with functor
compositions under strict definitional equality and demonstrate strategies for navigating dependent
type theory to construct fully formal, pointfree proofs of higher-level category theory results.
\end{abstract}


\section{Introduction}

The \emph{Co\!Yoneda Lemma} serves as a cornerstone of modern category theory, asserting that for
any category  $\mathcal{C}$, object $X \in \mathcal{C}$, and functor
$F : \mathcal{C} \to \mathbf{Set}$, natural transformations $\text{Nat}(\mathcal{C}(X, -), F)$ are
isomorphic to functor evaluations  $F(X)$. This result fundamentally states that an object $X$ is
completely determined by the morphisms out of it.

In computational settings, we frequently encounter categories equipped with additional structure.
Categories that have a functor mapping functions of types (or sets) to programs, their images being
called pure programs. We define such a category as a \emph{Functional Category}. In this paper, we
explore how a variation of the CoYoneda isomorphism extends in this setting. The variation replaces
elements by global elements and is entirely pointfree.

A key motivation for this exploration is the world of computer programs, which consist of both 
"pure programs without side-effects" and "impure programs with side-effects". We set out to
formulate a CoEndoYoneda equivalence for endofunctors that could model this broader programming
world using monads. Remarkably, we discover that the structural requirements of our equivalence
serve as a sieve for purity, mathematically forcing monadic referential transparency.

We provide a comprehensive discussion of how these concepts and theorems are formalized in the
Lean 4 proof assistant using the \texttt{Mathlib} library.

\section{Functional Categories}

Let $\mathcal{C}$ be a locally small category. We denote the category of types (or sets) as
$\mathbf{Type}$. For any object $X \in \mathcal{C}$, the covariant Hom-functor, which we will call
the CoYoneda functor, is defined as:
\[
CY\!F_X : \mathcal{C} \to \mathbf{Type}, \quad Y \mapsto \mathcal{C}(X, Y)
\]

\begin{definition}
A \textbf{Functional Category} is a category $\mathcal{C}$ equipped with the following data:
\begin{enumerate}
    \item A designated functor $\Phi : \mathbf{Type} \to \mathcal{C}$.
    \item A constructed endofunctor $GF : \mathcal{C} \to \mathcal{C}$ 
          defined by \\
          $GF(X) = \Phi(CY\!F_{\Phi(1)}(X))$.
    \item A monadic structure on the globals of $\mathcal{C}$, provided by a \\
          unit natural transformation  $\gamma\eta : \text{Id}_{\mathcal{C}} \to GF$ and a \\
          multiplication natural transformation $\gamma\mu : GF \circ GF \to GF$.
\end{enumerate}

These structural components must satisfy the standard monad left unit law: 
\[
\gamma\eta_{GF(X)} \circ \gamma\mu_X = \text{id}_{GF(X)}
\]
\end{definition}

For now no other extra monad laws are needed in the formal Lean 4 proof. This is an interesting
observation on its own.

This structure essentially allows $\mathcal{C}$ to reflect the category of types into itself, while
treating $GF$ as a monadic context. We define the \emph{Co\!Yoneda EndoFunctor}
$CY\!EF_X : \mathcal{C} \to \mathcal{C}$ for a fixed object $X$ as the composition:

\[
CY\!EF_X = \Phi \circ CY\!F_X
\]

We use the following definition of \emph{global element} as a morphisms from the canonical singleton
$\Phi(1)$ to $X$.

\[
G_{\mathit{\Phi}}(X) = \mathcal{C}(\Phi(1), X)
\]


\subsection{Formalization of Functional Categories in Lean 4}

In Lean 4, a functional category is modeled as a type class extending the standard
\texttt{LargeCategory} from \texttt{Mathlib}. 

\vspace{6pt}
\begin{lstlisting}[basicstyle=\tiny]
abbrev CYF {C : Type (u + 1)} [LargeCategory.{u} C] (X : C) : C ⥤ Type u :=
  coyoneda.obj (op X)

abbrev CYEF_def 
  {C : Type (u + 1)} [LargeCategory.{u} C] (Φ : Type u ⥤ C) (X : C) : C ⥤ C :=
    CYF X ⋙ Φ

abbrev GEF_def {C : Type (u + 1)} [LargeCategory.{u} C] (Φ : Type u ⥤ C) : C ⥤ C :=
  CYEF_def Φ (Φ.obj PUnit)

class FunctionalCategory (C : Type (u + 1)) extends LargeCategory.{u} C where
  Φ : Type u ⥤ C

  γμ : (GEF_def Φ ⋙ GEF_def Φ) ⟶ GEF_def Φ
  γη : 𝟭 _ ⟶ GEF_def Φ

  γ_left_unit : 
    ∀ (X : C), γη.app ((GEF_def Φ).obj X) ≫ γμ.app X = 𝟙 ((GEF_def Φ).obj X)

  φ {X Y : Type u} (f : X → Y) : Φ.obj X ⟶ Φ.obj Y := Φ.map (TypeCat.ofHom f)

  φ_eq : 
    ∀ {X Y : Type u} (f : X → Y), φ f = Φ.map (TypeCat.ofHom f) := 
      by intros; rfl

  γη_φ :
    ∀ (X : Type u), γη.app (Φ.obj X) = φ (fun x => φ (fun _ => x)) :=
      by intros; rfl
\end{lstlisting}
\vspace{6pt}

and global elements are modeled as

\vspace{6pt}
\begin{lstlisting}[basicstyle=\tiny]
abbrev G 
  {C : Type (u + 1)} [LargeCategory.{u} C] (Φ : Type u ⥤ C) (X : C) : Type u :=
    Φ.obj PUnit ⟶ X
\end{lstlisting}
\vspace{6pt}

A canonical, trivial instance of the type class is the world of pure programs without side-effects
(functions), $\mathbf{Type}$ itself, where $\Phi$ is the identity functor.

\vspace{6pt}
\begin{lstlisting}[basicstyle=\tiny]
instance typesFunctionalCategory : FunctionalCategory (Type u) where
  Φ := 𝟭 (Type u)
  γη := {
    app := fun X => (𝟭 (Type u)).map (↾(fun x => ↾(fun _ => x)))
    naturality := fun X Y f => by ext; rfl
  }
  γμ := {
    app := fun X => (𝟭 (Type u)).map (↾(fun (f : PUnit ⟶ PUnit ⟶ X) => ↾(fun pu => f pu pu)))
    naturality := fun X Y f => by ext; rfl
  }
  γ_left_unit := fun _ => by rfl
\end{lstlisting}
\vspace{6pt}

\section{The CoEndoYoneda Equivalence}

The requirements of \texttt{FunctionalCategory} prove to be sufficient to establish a CoEndoYoneda
equivalence for endofunctors.

\subsection{Auxiliary definitions and theorems}

We construct a designated global element from a designated natural transformation as follows 

\vspace{6pt}
\begin{lstlisting}[basicstyle=\tiny]
def τX2ggfx {F : C ⥤ C} {X : C} (τX : CYEF X ⟶ F ⋙ GEF) :
  Global ((F ⋙ GEF).obj X) :=
    let ggfx := φ (fun _ => 𝟙 X) ≫ τX.app X
    ggfx
\end{lstlisting}
\vspace{6pt}

We construct a designated (natural, not shown) transformation from a designated global element as
follows 

\vspace{6pt}
\begin{lstlisting}[basicstyle=\tiny]
def ggfx2τX
  {F : C ⥤ C} {X : C} (ggfx : Global (GEF.obj (F.obj X))) : CYEF X ⟶ (F ⋙ GEF) :=
    let τX : CYEF X ⟶ F ⋙ GEF := {
      app Y := φ (ggfx ≫ (F ⋙ GEF).map .) ≫ γμ.app (F.obj Y)
      naturality _ _ h := ...
    }
    τX
\end{lstlisting}
\vspace{6pt}

We formulate (and prove, not shown) a left inverse theorem as follows

\vspace{6pt}
\begin{lstlisting}[basicstyle=\tiny]
theorem left_inverse {F : C ⥤ C} {X : C} (τX : CYEF X ⟶ (F ⋙ GEF)) :
    τx = ggfx2τx (τx2ggfx τx) := by ...
\end{lstlisting}

We formulate (and prove, not shown) a right inverse theorem as follows

\vspace{6pt}
\begin{lstlisting}[basicstyle=\tiny]
theorem right_inverse {F : C ⥤ C} {X : C} (ggfx : Global ((F ⋙ GEF).obj X)) :
    ggfx = τx2ggfx (ggfx2τx ggfx) := by ...
\end{lstlisting}
\vspace{6pt}

Notice that both inverses apply $\varphi$ (defined as $\Phi.\text{map}$) to 
$\text{fun } \_ \Rightarrow \text{id}_X$, which mirrors the classic CoYoneda construction.

\subsection{CoEndoYoneda Equivalence}

Utilizing the auxiliary definitions and theorems above we can now define
the main equivalence.

\vspace{6pt}
\begin{lstlisting}[basicstyle=\tiny]
def coEndoYonedaEquiv {F : C ⥤ C} {X : C} : 
  (CYEF X ⟶ (F ⋙ GEF)) ≃ Global ((F ⋙ GEF).obj X) where
  toFun     := τX2ggfx
  invFun    := ggfx2τX
  left_inv  := fun τX => (left_inverse  τX).symm
  right_inv := fun τX => (right_inverse τX).symm
\end{lstlisting}
\vspace{6pt}

This equivalence provides a completely pointfree and logically verified method for evaluating
and constructing natural transformations in functional categories.

\section{The Essence of Purity \\ (Every Disadvantage Has an Advantage)}

The original motivation was to define a CoEndoYoneda equivalence that could play a role in the
broader programming world, which consists of both pure programs (functions) and impure programs
(computation with side-effects valued functions). The natural setting for such impure programs is
the Kleisli category of a Monad $M$, denoted \texttt{KleisliCat M}.

However, attempting to instantiate \texttt{KleisliCat M} as a \texttt{FunctionalCategory} reveals
an apparent limitation.  The equivalence cannot be instantiated for arbitrary monads.  Thus, the
\texttt{coEndoYonedaEquiv} is not universally useful for the (potentially impure) programming
world.

But this disadvantage reveals a profound mathematical advantage: \\
the \texttt{FunctionalCategory} requirements, more precisely $\gamma\eta$-naturality, act as a
strict classification of purity for \texttt{KleisliCat M}. They mathematically declare the essence
of referential transparency of monad computations
(cfr. referential transparency of expression evaluations).

\subsection{The Purity Constraint}

The core constraint imposed by $\gamma\eta.\text{naturality}$ demands that mapping \texttt{pure}
over the structure of a monadic value $m$ is exactly the same as wrapping the entire structure $m$
in \texttt{pure}. Formally, it forces the equality:
\[
m \gg= \lambda z. \text{pure}(\text{pure}(z)) = \text{pure}(m)
\]

If a Monad satisfies this purity constraint (which is strictly required to instantiate
\texttt{FunctionalCategory} via \texttt{KleisliCat}), then its shape functor evaluated at
\texttt{PUnit} is a subsingleton (proof not shown). 

\vspace{6pt}
\begin{lstlisting}[basicstyle=\tiny]
def EnforcesPurity : Prop :=
  ∀ {X : Type v} (mx : M X), (mx >>= fun x => pure (pure x)) = pure mx

theorem purity_implies_subsingleton_shape
  (h_pure : EnforcesPurity M) : Subsingleton (M PUnit) := by ...
\end{lstlisting}
\vspace{6pt}

This implies that $M$ is a "Subsingleton Monad" 
(such as the Identity Monad or the Terminal Monad). 
In other words, any standard computational effect 
(like State, IO, or Option) is filtered out by the type system. 

We formalize this conclusion explicitly (proof not shown):
\vspace{6pt}
\begin{lstlisting}[basicstyle=\tiny]
theorem functional_category_implies_purity
  (γη : 𝟭 (KleisliCat M) ⟶ (GEF_def (κΦ M)))
  (h_γη_app : ∀ X, γη.app X = fun x => pure (fun _ => pure x)) :
  EnforcesPurity M := by ...
\end{lstlisting}
\vspace{6pt}

{\em Referential transparency is exactly the requirement for a computational context to admit a
global coyoneda endofunctor equivalence}.

\section{Proof Techniques and Challenges in Lean 4}

Formalizing abstract category theory in dependent type theory (DTT) introduces unique
difficulties, particularly concerning definitional equality. 

\subsection{Strict Functor Composition}
In standard mathematics, functor composition is strictly associative. In Lean 4, however, these are
represented as distinct abstract syntax trees. They are \emph{isomorphic} via 
\texttt{Category.assoc}, but they are not \emph{definitionally equal}. 

When writing the proof of \texttt{left\_inverse}, this limitation prevents standard
term-rewriting tactics like \texttt{simp} from effortlessly reducing deeply nested compositions.

\subsection{Constructing Equality Chains}
To circumvent the strict type-checker without resorting to manual casts (such as \texttt{eqToHom}),
the proof for \texttt{left\_inverse} abandons the monolithic \texttt{simp} approach in favor of
explicit equality chains acting on the natural transformation components.

By utilizing \texttt{Eq.trans} and \texttt{congr\_arg}, we explicitly guide the type checker
through the unit laws and naturality squares:

\begin{lstlisting}[basicstyle=\tiny]
  have h1 : 
    τX.app Y = τX.app Y ≫ 𝟙 ((GEF_def Φ).obj (F.obj Y)) := 
      (Category.comp_id (τX.app Y)).symm
  have h2 : 
    𝟙 ((GEF_def Φ).obj (F.obj Y)) = 
      γη.app ((GEF_def Φ).obj (F.obj Y)) ≫ γμ.app (F.obj Y) :=
        (γ_left_unit (F.obj Y)).symm
  have h_left_unit : 
    τX.app Y = 
      τX.app Y ≫ (γη.app ((GEF_def Φ).obj (F.obj Y)) ≫ γμ.app (F.obj Y)) :=
        Eq.trans h1 (congr_arg (τX.app Y ≫ .) h2)
\end{lstlisting}
\vspace{6pt}

This localized approach bypasses the rigid typing of the category morphisms by
acting directly on the equality type \texttt{=} in the universe of types. This technique is highly
robust and avoids the infamous "motive is not type correct" errors common in Lean category theory
formalization.

\section{Conclusion}

The Pointfree CoEndoYoneda Lemma is a variation of the classical Pointful CoYoneda Lemma in a
functional category setting. While our initial aim was to generalize this equivalence to arbitrary
computational effects via monads, the rigorous formalization process revealed that the required
categorical structure fundamentally filters out impurities.

By successfully formalizing the lemma in Lean 4, we have not only verified its absolute logical
correctness but also established that \texttt{FunctionalCategory} mathematically classifies
referential transparency. Moreover, we have demonstrated a concrete methodology for handling strict
composition equality in dependent type theory. The approach of combining explicit structural
definitions with targeted, step-by-step equality chaining provides a reliable method for
formalizing advanced category theory in modern proof assistants.

\end{document}
